
\section[Aplicación a los recubrimientos principales]{Aplicación a los
recubrimientos principales: teorías de Kummer y de Artin--Schreier}
\label{XI.6}


\begin{proposition}
\label{XI.6.1}
Sean~$S$ un preesquema y~$n$ un entero~$>0$; sea
$u_n\colon\GG_{m\,S}\to\GG_{m\,S}$ el homomorfismo de elevación a la
potencia $n$-ésima, y sea~$\bbmu_{n\,S}$
\label{indnot:kb}\oldindexnot{$\bbmu_{n\,S}$|hyperpage}%
su núcleo. Entonces~$\bbmu_{n\,S}$ es finito y localmente libre de rango~$n$
sobre~$S$, y es étale sobre~$S$ si y solo si, para todo~$s\in S$, la
característica de~$s$ es prima con~$n$. La sucesión de homomorfismos
$$
0\to \bbmu_{n\,S}\to\GG_{m\,S}\lto{u_n}\GG_{m\,S}\to 0
$$
es exacta en el sentido de~\Ref{XI.4}. (La llamaremos la \emph{sucesión
exacta de Kummer}
\index{Kummer (sucesión exacta de)|hyperpage}%
sobre~$S$, relativa al entero~$n$).
\end{proposition}

Se tiene
\marginpar{303}
$$
\GG_m=\Spec\cal{O}_S[t,t^{-1}],
$$
y~$u_n$ corresponde al homomorfismo~$u_n$ de las
$\cal{O}_S$-álgebras afines, dado por
$$
u_n(t)=t^n,
$$
por otra parte, la sección unidad de~$\GG_{m\,S}$ corresponde al
homomorfismo de aumento de $\cal{O}_S$-álgebras dado por
$$
\varepsilon(t)=1,
$$
cuyo núcleo es, por tanto, el
\ifthenelse{\boolean{orig}}
{ideal}
{ideal}
principal~$(t-1)$. La imagen de este último por~$u_n$ es, pues, el ideal
principal~$(1-t^n)$, y se obtiene
$$
\bbmu_{n\,S}=\Spec\cal{O}_S[t]/(1-t^n),
$$
lo que muestra, en particular, que~$\bbmu_{n\,S}$ es finito sobre~$S$ y está
definido por un álgebra sobre~$S$ que es libre de rango~$n$, con base formada
por los~$t^i$ ($0\leq i\leq n-1$). Para que sea étale en~$s\in S$, es
necesario y suficiente que el álgebra reducida~$k[t]/(1-t^n)$, donde
$k=\kres(s)$, obtenida adjuntando formalmente las raíces $n$-ésimas de la
unidad a~$k$, sea separable sobre~$k$, \,\ie que las raíces de~$1-t^n$ en una
clausura algebraica de~$k$ sean todas distintas, lo cual equivale a que~$n$
sea primo con la característica. Por último, para demostrar que la sucesión
de homomorfismos de~\Ref{XI.6.1} es exacta, en virtud del
criterio~\Ref{XI.4.2} se reduce uno a probar que
\ifthenelse{\boolean{orig}}
{$v$}
{$u_n$}
es fielmente plano. Podemos suponer, evidentemente, que~$S$ es afín de
anillo~$A$, de modo que~$\GG_{m\,S}$ es afín de anillo
$B=A[t,t^{-1}]$; basta comprobar que~$u_n$ hace de~$B$ un módulo libre de
rango~$n$ sobre~$B$, o, lo que es lo mismo, que~$u_n$ es inyectivo y que
$A[t,t^{-1}]$ es un módulo libre de rango~$n$ sobre~$A[t^n,t^{-n}]$. En
efecto, se comprueba fácilmente que los~$t^i$ ($0\leq i\leq n-1$) forman
una base del primero sobre el segundo, lo que termina la demostración.

\begin{definition}
\label{XI.6.2}
Se
\marginpar{304}
llama a~$\bbmu_{n\,S}$ el \emph{grupo de Kummer de rango~$n$ sobre~$S$}, y
se llama \emph{recubrimiento principal de Kummer de rango
\ifthenelse{\boolean{orig}}
{$n\,S$}
{$n$ sobre~$S$}}
a todo fibrado homogéneo principal sobre~$S$ cuyo grupo es el grupo de
Kummer de rango~$n$.
\end{definition}

El conjunto de estos recubrimientos es un grupo, denotado por
$\H^1(S,\bbmu_{n\,S})$ o, simplemente, por~$\H^1(S,\bbmu_n)$. Obsérvese que
la formación del grupo de Kummer de rango~$n$ sobre~$S$ es compatible con la
extensión de la base, de modo que~$\bbmu_{n\,S}$ proviene por extensión de la
base del \emph{grupo de Kummer absoluto~$\bbmu_n$} sobre~$\Spec(\ZZ)$.

Denotemos por~$(\ZZ/n\ZZ)_S$ el $S$-grupo definido por el grupo finito
ordinario~$\ZZ/n\ZZ$. Si~$G$ es un $S$-grupo cualquiera, los homomorfismos
de $S$-grupos~$u$ de~$(\ZZ/n\ZZ)_S$ en~$G$ corresponden biyectivamente, y
de manera compatible con el cambio de base, a las secciones de~$G$ sobre~$S$
cuya potencia $n$-ésima es la sección unidad, haciendo corresponder a~$u$ la
imagen por~$u$ de la sección de~$(\ZZ/n\ZZ)_S$ sobre~$S$ definida por el
generador~$1\bmod n\ZZ$ de~$\ZZ/n\ZZ$. Dicho esto:

\begin{corollaire}
\label{XI.6.3}
Si~$\bbmu_{n\,S}$ es étale sobre~$S$, se obtiene así una correspondencia
biyectiva entre los isomorfismos de $S$-grupos
$(\ZZ/n\ZZ)_S\isomto\bbmu_{n\,S}$ y las secciones de~$\cal{O}_S$ que son
de orden exactamente~$n$ en cada componente conexa de~$S$ (una sección tal
se llamará «raíz primitiva $n$-ésima de la unidad sobre~$S$»). Por tanto,
para que~$\bbmu_{n\,S}$ sea isomorfo, como $S$-grupo, a~$(\ZZ/n\ZZ)_S$, es
necesario y suficiente que sea étale
\ifthenelse{\boolean{orig}}
{sobre~$S$}
{sobre~$S$,}
\ie que las características residuales de~$S$ sean primas con~$n$, y que
exista una raíz primitiva $n$-ésima de la unidad sobre~$S$.
\end{corollaire}

Esto explica el papel que desempeñan en la teoría clásica de Kummer la
hipótesis de que el cuerpo base (que desempeña el papel de~$S$) tenga
característica prima con~$n$ y contenga las raíces $n$-ésimas de la unidad,
y la elección de una raíz primitiva $n$-ésima de la unidad. Una vez que se
dispone del lenguaje de los esquemas, ya no hay por qué imponerse estas
hipótesis, y conviene razonar directamente sobre~$\bbmu_n$ en vez
de~$\ZZ/n\ZZ$. Así, la conjunción de~\Ref{XI.6.1}, \Ref{XI.4.5}
y~\Ref{XI.5.3} nos da la relación general siguiente entre la teoría de los
recubrimientos principales de Kummer y la de los grupos de Picard:

\begin{proposition}
\label{XI.6.4}
Sea
\marginpar{305}
$S$ un preesquema. Se tiene una sucesión exacta canónica
$$
\textstyle
0\to\H^0(S,\bbmu_n)\to\H^0(S,\cal{O}^*_S)\to
\H^0(S,\cal{O}^*_S)\to\H^1(S,\bbmu_n)\to
\H^1(S,\cal{O}^*_S)\to\H^1(S,\cal{O}^*_S),
$$
de donde, poniendo
\ifthenelse{\boolean{orig}}
{$\H^1(S,\cal{O}_S)=\Pic(S)$,}
{$\H^1(S,\cal{O}_S^*)=\Pic(S)$,}
y denotando, para todo grupo abeliano~$A$, por~${_n}A$ y~$A_n$ el núcleo y
el conúcleo de la multiplicación por~$n$ en~$A$, se obtiene la sucesión exacta
$$
0\to
\ifthenelse{\boolean{orig}}
{\H^0(S,\cal{O}_S)^*_n}
{\H^0(S,\cal{O}_S^*)_n}
\to\H^1(S,\bbmu_n)\to
{\sideset{_n}{}\Pic}(S)\to 0.
$$
\end{proposition}

Vamos a explicitar dos casos importantes en los que uno u otro término
extremo de esta sucesión exacta es nulo:

\begin{corollaire}
\label{XI.6.5}
Supongamos que~$\sideset{_n}{}\Pic(S)=0$ (por ejemplo, que~$S$ sea el
espectro de un anillo local o de un anillo factorial), y sea~$A$ el anillo
$\H^0(S,\cal{O}_S)$. Entonces se tiene un isomorfismo canónico
$$
\H^1(S,\bbmu_n)\isomto A^*/{A^*}^n.
$$
\end{corollaire}

Este es esencialmente el enunciado clásico de la teoría de Kummer cuando~$S$
es el espectro de un cuerpo.

\begin{corollaire}
\label{XI.6.6}
Supongamos que todo elemento de~$\H^0(S,\cal{O}_S)$ es una potencia
$n$-ésima; por ejemplo, que~$\H^0(S,\cal{O}_S)$ sea un
\ifthenelse{\boolean{orig}}
{compuesto}
{producto}
de cuerpos algebraicamente cerrados, o que~$S$ sea reducido y propio sobre
un cuerpo algebraicamente cerrado~$k$. Entonces se tiene un isomorfismo
canónico
$$
\H^1(S,\bbmu_n)\isomto\sideset{_n}{}\Pic(S).
$$
\end{corollaire}

En particular, cuando~$S$ es propio y conexo sobre un cuerpo algebraicamente
cerrado~$k$, esto relaciona el grupo fundamental de~$S$ con los puntos de
orden finito del esquema de Picard~$P$ de~$S$ sobre~$k$; así, para~$n$ primo
con la característica, se tiene un isomorfismo
$$
\Hom(\pi_1(S),\ZZ/n\ZZ)\simeq {_n P}(k)
$$
que se utiliza a menudo en geometría
\marginpar{306}
algebraica. Como aplicación, cuando la componente conexa~$P^0$ de~$P$ es un
esquema en grupos completo, de dimensión~$g$, se ve, utilizando los
resultados recordados en~\Ref{XI.2} y la finitud del grupo de torsión de
\ifthenelse{\boolean{orig}}
{Néron--Sévéri}
{Néron--Severi}
que, para todo número primo~$\ell$ distinto de la característica, la
componente $\ell$-primaria del grupo fundamental abelianizado~$\pi_1(S)$ es
un módulo de tipo finito y de rango~$2g$ sobre el anillo~$\ZZ_\ell$ de los
enteros $\ell$-ádicos (y, además, libre excepto, a lo sumo, para un número
finito de valores de~$\ell$). Como observó Serre, esto permite demostrar,
bajo ciertas condiciones, que cuando~$X$ es un esquema plano y proyectivo
sobre~$S$ conexo, los esquemas de Picard de las fibras de~$X$ tienen todos
la misma dimensión, aplicando el teorema de semicontinuidad
(\Ref{X.2.3}); el argumento de Serre se aplica siempre que exista el esquema
de Picard de~$X$ sobre~$S$ y los esquemas de Picard conexos de las fibras
de~$X$ sobre~$S$ sean esquemas en grupos propios, por ejemplo cuando las
fibras geométricas de~$X$ sobre~$S$ sean normales (manteniéndose~$X$ plano y proyectivo
sobre~$S$), en particular si~$X$ es liso y proyectivo sobre~$S$.

Sea ahora~$p$ un número primo, y supongamos que~$S$ es un preesquema de
característica~$p$, \,\ie tal que~$p\cdot\cal{O}_S=0$. Entonces el
homomorfismo de elevación a la potencia $p$-ésima en~$\cal{O}_S$ es aditivo,
y el morfismo correspondiente, obtenido reemplazando~$S$ por un~$T$ variable
sobre~$S$,
$$
F\colon\GG_{a\,S}\to\GG_{a\,S}
$$
es, por tanto, un homomorfismo de $S$-grupos, llamado el \emph{homomorfismo
de Frobenius}. (N.B. Un morfismo tal está definido para todo $S$-preesquema
$G$ que provenga por extensión de la base de un preesquema~$G_0$ sobre el
cuerpo primo~$\ZZ/p\ZZ$, y este morfismo es un homomorfismo de grupos si~$G_0$
es un preesquema en grupos). Pondremos
$$
\wp=\id-F\colon\GG_{a\,S}\to\GG_{a\,S}.
$$
